\section{Introduction}
 
In recent years, Large Language Models (LLMs) such as GPT \cite{brown2020language} have demonstrated remarkable capabilities in modeling complex, high-dimensional data across various domains, including text and image generation \cite{zhao2023survey,liu2024visual}. Inspired by these successes, there has been significant interest \cite{agostinelli2023musiclm,borsos2023audiolm,wang2023neural,yang2023uniaudio} in exploring the application of LLMs to audio generation.

Audio codecs \cite{zeghidour2021soundstream} have emerged as a critical technique for audio LLMs, bridging the gap between continuous audio waveforms and token-based language models. By discretizing high-rate audio signals into a finite set of tokens, these codecs enable the application of LLM architectures to audio data, leveraging the successes of textual LLMs. 

However, prior research on audio codecs has primarily focused on achieving lower compression rates and higher reconstruction quality \cite{kumar2024high,defossez2022high,yang2023hifi}. Meanwhile, many efforts in audio generation have concentrated on enhancing model architecture, scaling, or leveraging larger datasets. For instance, AudioLM \cite{borsos2023audiolm} adopts a two-stage pipeline that models the acoustic token in an autoregressive way conditioned on the semantic token. VALL-E \cite{wang2023neural}, the first TTS framework to leverage large, diverse, and multi-speaker speech data, demonstrates strong in-context learning capabilities similar to GPT-3, treating TTS as a language modeling task on audio codecs. MusicGen \cite{copet2024simple} generates music using a single-stage transformer LM alongside efficient token interleaving patterns. Similarly, UniAudio \cite{yang2023uniaudio} scaled up to 165K hours of audio and 1B parameters, utilizing LLM techniques to generate tokens for various types of audio, including speech, sounds, music, and singing, given different input conditions.

While these works have shown success in developing audio language models, they all rely on the acoustic codecs such as Encodec \cite{defossez2022high} or Soundstream \cite{zeghidour2021soundstream} for audio tokenization and de-tokenization. However, these acoustic codecs were originally designed for audio compression rather than for audio language models. This misalignment means the design may not be optimal for audio language modeling.

To design a better audio codec for Audio LLMs, we drew inspiration from the initial purpose of LLMs such as GPT, which were designed to process text. These models focus on understanding and generating natural language, which is inherently rich in semantics. Motivated by this, we assume that a better audio tokenizer should encapsulate rich semantic information to facilitate an easy understanding of audio content, thus reducing the language model's burden in interpreting tokens. However, most audio codecs focus on acoustic reconstruction which ignores the semantic information. As a result, LLM essentially tries to predict the local fluctuations of the audio signal, which is difficult, and methods like VALL-E, which condition acoustic token generation on text transcriptions, frequently result in content inaccuracies causing elevated word error rates (WER), stemming from the semantic misinterpretations of acoustic tokens, leading to word skipping and errors.

To address this issue, approaches like SpeechTokenizer \cite{zhang2023speechtokenizer} have attempted to disentangle speech into separate tokens for content and timbre and perform distillation-based semantic and acoustic integration. However, this method may not integrate smoothly with all audio LLMs, especially those requiring uniform token treatment across different layers, such as utilizing flattened codec tokens \cite{yang2023uniaudio,copet2024simple}. 

In this paper, We propose a straightforward yet effective method termed ``X-codec'', which integrates both semantic and acoustic features into a unified tokenization framework. The X-Codec architecture employs a distinctive ``X-shaped'' structure, characterized by two inputs and two outputs, unifying semantic and acoustic information within a single Residual Vector Quantizer (RVQ) structure. This design enables simultaneous embedding learning of semantic richness and acoustic fidelity for every token, resulting in better performance for audio LLM. 

We have conducted comprehensive evaluations of X-Codec across various applications, including text-to-speech, music continuation, and text-to-sound synthesis. The results consistently demonstrate the effectiveness of the proposed method. Furthermore, our comparative evaluation on VALL-E based TTS demonstrates that X-Codec outperforms existing disentanglement techniques, thereby highlighting its efficacy and versatility in advancing audio LLM technologies.






% In recent years, Large Language Models (LLMs) such as GPT \cite{brown2020language} have demonstrated remarkable capabilities in modeling complex, high-dimensional data across various domains, including text and image generation \cite{zhao2023survey,liu2024visual}. Inspired by these successes, there has been significant interest \cite{agostinelli2023musiclm,borsos2023audiolm,wang2023neural,yang2023uniaudio} in exploring the application of LLMs to audio generation.

% Audio codecs \cite{zeghidour2021soundstream} have emerged as a critical technique for audio LLMs, bridging the gap between continuous audio waveforms and token-based language models. By discretizing high-rate audio signals into a finite set of tokens, these codecs enable the application of LLM architectures to audio data, leveraging the successes of textual LLMs. 

% However, prior research on audio codecs has primarily focused on achieving lower compression rates and higher reconstruction quality \cite{kumar2024high,defossez2022high,yang2023hifi}. Meanwhile, many efforts in audio generation have concentrated on enhancing model architecture, scaling, or leveraging larger datasets. For instance, AudioLM \cite{borsos2023audiolm} adopts a two-stage pipeline that models the acoustic token in an autoregressive way condition on the semantic token. VALL-E \cite{wang2023neural}, the first TTS framework to leverage large, diverse, and multi-speaker speech data, demonstrates strong in-context learning capabilities similar to GPT-3, treating TTS as a language modeling task using audio codecs. MusicGen \cite{copet2024simple} generates music using a single-stage transformer LM alongside efficient token interleaving patterns. Similarly, UniAudio \cite{yang2023uniaudio} scaled up to 165K hours of audio and 1B parameters, utilizing LLM techniques to generate tokens for various types of audio, including speech, sounds, music, and singing, based on given input conditions.

% While these works have shown success in developing audio language models, they all rely on the   Encodec \cite{defossez2022high} or Soundstream \cite{zeghidour2021soundstream}for audio tokenization and de-tokenization. However, these codecs were originally designed for audio compression rather than for audio language models. This misalignment means the design may not be optimal for audio language modeling.


% % 放到related work
% % There have been some attempts to design better tokenizers for audio language models. For example, SpeechTokenizer adopts a disentanglement approach, using HuBERT to separate speech into the first VQ representing content and the remaining VQ representing timbre and acoustic details. This allows VALL-E's autoregressive (AR) stage to better model the content’s context, while the non-autoregressive (NAR) stage can supplement acoustic details. However, this disentanglement approach might be more suited to architectures like VALL-E and may not be as effective for LLMs that aim to treat each token fairly, such as when flattening the tokens of codec. Additionally, the assumption of separating content and timbre is challenging to generalize to other types of audio, such as music and sound. 

% To design a better audio codec for Audio LLMs, we drew inspiration from the initial purpose of LLMs such as GPT, which were designed to process text. These models focus on understanding and generating natural language, which is inherently rich in semantics. Motivated by this, we assume that a better audio tokenizer should encapsulate rich semantic information to facilitate an easy understanding of audio content, thus reducing the language model's burden in interpreting tokens. However, most audio codecs focus on acoustic reconstruction which ignores the semantic information. As a result, methods like VALL-E, which condition acoustic token generation on text transcriptions, frequently result in content inaccuracies causing elevated word error rates (WER), stemming from the semantic misinterpretations of acoustic tokens, leading to word skipping and errors.

% To address this issue, some approaches like SpeechTokenizer have attempted to disentangle speech into distinct tokens for content and timbre. Nevertheless, this method may not seamlessly align with all audio LLM architectures, particularly those requiring a uniform treatment of different layer tokens, such as approaches that utilize flattened codec tokens\cite{yang2023uniaudio,copet2024simple}.  
% To bridge this gap, We propose a straightforward but effective method termed X-codec, which innovative integrates both semantic and acoustic features into a unified tokenization framework. The X-Codec architecture employs a distinctive ‘X-shaped' structure, characterized by two inputs and two outputs, unifying semantic and acoustic information within a single Residual Vector Quantizer (RVQ) system. This structural design enables simultaneous embedding learning of semantic richness and acoustic fidelity for every token, resulting in better performance for audio LLM. We  evaluated X-Codec through extensive experiments across applications in text-to-speech, music continuation, and text-to-sound, validating our approach. Evaluation on VALL-E further confirms that X-Codec surpasses pervious disentanglement methods in performance, underscoring its efficacy and versatility in advancing audio LLM technologies.


% % There are primarily two types of audio tokens \todo{why good/bad}, as identified in recent research \cite{borsos2023audiolm}: acoustic and semantic. Acoustic tokens focus on compressing audio to achieve perceptual equivalence, including SoundStream \cite{zeghidour2021soundstream}, Encodec \cite{defossez2022high}, and DAC \cite{kumar2024high}. Conversely, semantic tokens prioritize retaining the content's semantic integrity through techniques like clustering in self-supervised learning models, such as HuBERT \cite{hsu2021hubert} and Wav2Vec-BERT \cite{barrault2023seamless}. 

% % While acoustic tokens typically allow for high-fidelity audio generation, they may lack in phoneme distinction, making them less suitable for capturing linguistic content. Semantic tokens, in contrast, provide long-term coherence but often suffer from poorer reconstruction quality. 
% % Therefore, two-stage generation process is commonly adopted, semantic tokens are first extracted and then transformed into acoustic tokens with coarse-to-fine modeling, which are used to generate the waveform afterwards. 
% % Typical systems include SPEAR-TTS \cite{kharitonov2023speak} and MusicLM \cite{agostinelli2023musiclm}, for speech and music generation, respectively. Such two-stage process is complicated, suboptimal since semantic and acoustic tokens are separately developed, and introduces error accumulation. To truely perform language-like modelling on audio, an unified audio codec is essentially needed to comprehensively supports both high-fidelity audio reconstruction and robust semantic understanding in a single-stage, decoder-only generation process.  

% % To build a unified audio tokenizer, several properties must be met:
% % 1) The tokenizer should encapsulate rich semantic information to facilitate an easy understanding of audio content, thus reducing the generative model's burden in interpreting tokens.
% % 2) It should enable the reconstruction of audio with high fidelity.
% % 3) It should support straightforward, one-stage GPT-style audio generation using a single unified tokenizer.
% % 4) It should be versatile enough to handle all types of audio, including speech, music, and environmental sounds.

 
% % To address the inherent limitations of the traditional two-stage audio processing pipeline, this paper introduces a novel audio tokenizer, X-Codec, which innovative integrates both semantic and acoustic features into a unified tokenization framework. The X-Codec architecture employs a distinctive ‘X-shaped' structure, characterized by two inputs and two outputs, unifying semantic and acoustic information within a single Residual Vector Quantizer (RVQ) system. This structural design enables simultaneous embedding learning of semantic richness and acoustic fidelity, allowing the codec to capture a comprehensive audio profile cohesively.
% % We evaluate X-Codec's performance through extensive experiments on semantic fidelity and reconstruction quality. Additionally, we explore its application in a GPT-style, decoder-only model tailored for audio generation. 

% % Our contributions are threefold:
% % 1) We introduce the X-Codec, a unified audio tokenization approach that elegantly integrates semantic and acoustic information.
% % 2) We present empirical studies that demonstrate X-Codec's superior capabilities in semantic preservation and audio quality.
% % 3) We explore the application potentials of X-Codec in GPT-style audio generation, showing promising generation results.


  